\documentclass[monospace]{sfs-2487-2024}

% Muistio, joka on laadittu tasavälisellä kirjasinlajilla (Courier).
% Vaihtoehto [monospace] latoo asiakirjan kirjoituskonetyylisesti.
% Mukana sisäkkäinen luettelo, tasavälinen koodimerkintä ja koodilohko,
% jotka vastaavat markdown-version vastaavia merkintöjä.

\logo{\includegraphics{logo-organisaatio}}
\doctype{Muistio}
\date{5.6.2024}
\author{Tero Testi}

\subject{Tietojärjestelmä}
\title{Käyttöoikeuksien hallinta}

\contactinfo{Organisaatio Oy}{%
  Katuosoite\\
  Postinumero Postitoimipaikka\\
  it@organisaatio.fi}

\begin{document}

\maketitle

\marginlabel{Vastaanottajat}

Järjestelmänvalvojat

\section{Nykytilanne}

Käyttöoikeuksia hallitaan tällä hetkellä manuaalisesti jokaisen
järjestelmän omassa hallintarajapinnassa. Tunnuksia luodaan
pyydettäessä sähköpostilla, eikä vanhentuneita tunnuksia poisteta
järjestelmällisesti.

\section{Ehdotus}

Otetaan käyttöön keskitetty identiteetinhallintajärjestelmä, joka
tukee \texttt{SCIM}- ja \texttt{LDAP}-rajapintoja. Se yhdistää kaikki
organisaation järjestelmät yhteen hallintapisteeseen ja mahdollistaa
automaattisen tunnusten elinkaarihallinnan.

\begin{enumerate}
  \item Valitaan järjestelmä ja kilpailutetaan toimittajat.
  \item Toteutetaan pilotti kahdelle järjestelmälle:
    \begin{itemize}
      \item sähköpostijärjestelmä
      \item asianhallintajärjestelmä
    \end{itemize}
  \item Laajennetaan käyttöönotto koko organisaatioon.
\end{enumerate}

Tunnusten synkronointi ajastetaan komennolla:

\begin{verbatim}
idm sync --scope all --remove-stale
\end{verbatim}

\forinformation{Johtoryhmä}

\makecontactinfo

\end{document}
